[V.] Art.\ $[7.]$ Es bezeichnet hier $\mathrm{Dec.}\frac mM$ wie in Art.\ $1$ des vorhergehenden Bruchstücks [IV] den Werth von
\[
\Sigma(-1)^X\Sigma(-1)^Y\mathrm{Int.}\,p
\]
worin die Summation über alle ganze Zahlen $X$ und $Y$ auszudehnen ist, für welche die zugehörigen $\xi$ und $\eta$ innerhalb der Grenzen $0$ und $\tfrac12$ liegen.

Die Formeln für den Decidenten in Art.\ $7$ und $15$ sind nach der Angabe des Textes auf zwei be%
\SourceLeaf{388}%
sondern Wegen gefunden, um aber diese Erläuterungen nicht zu sehr auszudehnen, werden sie hier aus einer gemeinsamen Quelle abgeleitet.

Indem $X$ irgend einen bestimmten ganzzahligen Werth annimmt, sei $Y^{*}$ das kleinere, $Y^{**}$ das grössere der beiden $Y$, welche den Grenzwerthen von $\xi$, $\eta$ entsprechen. Die zu $Y^{*}$ und $Y^{**}$ zugehörigen Werthe von $p$ seien $p^{*}$ und $p^{**}$, die ebenso wie $Y^{*}$ und $Y^{**}$ einander nicht gleich werden können, weil die Summe $\Sigma$ sich nicht über die Grenzwerthe von $\xi$ und $\eta$ erstreckt.

Führt man auf dieselbe Weise wie in den beiden vorhergehenden Aufsätzen [III] und [IV] die Summation über alle bei demselben $X$ Statt habenden Werthe von $Y$ aus, setzt dabei für die Anzahl der zwischen $Y'$ und $Y''$ liegenden ungeraden Zahlen $[\tfrac12Y''-\tfrac12]-[\tfrac12Y'-\tfrac12]$ und fügt die Intensoren, die sich auf die Grenzen $\xi=0$ und $=\tfrac12$ beziehen, zwei Mal aber mit entgegengesetzten Zeichen hinzu, so erhält man für $\Sigma(-1)^Y\mathrm{Int.}\,p$ den aus sieben Theilen bestehenden Ausdruck
\[
\begin{gathered}
-\Sigma[-\mathrm{Int.}(p-\mu\omega i)+\mathrm{Int.}(p+\mu\omega i)]\\
\text{worin alle }p\text{ aufzunehmen, für welche }[Y]\text{ gerade,}\\
x\text{ oder }y\text{ ganz, incl. }\xi=0\text{ und }\tfrac12,\quad
\text{excl. }\eta=0\text{ und }\tfrac12
\end{gathered}
\]
\[
\left.\begin{aligned}
&-\mathrm{Int.}(p^{*}-\mu\omega i)\text{ wenn }[Y^{*}]\text{ gerade}\\
&+\mathrm{Int.}(p^{**}+\mu\omega i)\text{ wenn }[Y^{**}]\text{ gerade}
\end{aligned}\right\}\ \xi=0\text{ oder }\tfrac12,\quad0<\eta<\tfrac12
\]
\[
\left.\begin{aligned}
&-\mathrm{Int.}(p^{*}+\mu\omega i)\text{ wenn }[Y^{*}]\text{ gerade}\\
&+\mathrm{Int.}(p^{**}-\mu\omega i)\text{ wenn }[Y^{**}]\text{ gerade}
\end{aligned}\right\}\ 0<\xi<\tfrac12,\quad\eta=0\text{ oder }\tfrac12
\]
\[
\left.\begin{aligned}
&-\mathrm{Int.}(p^{*}+\mu\omega i)\text{ wenn }[Y^{*}]\text{ gerade}\\
&+\mathrm{Int.}(p^{**}-\mu\omega i)\text{ wenn }[Y^{**}-\omega]\text{ gerade}
\end{aligned}\right\}\ \xi=0\text{ oder }\tfrac12,\quad\eta=0\text{ oder }\tfrac12
\]
welcher mit $(-1)^X$ multiplicirt und über alle ganzzahligen $X$ summirt den Decidenten $\frac mM$ ergiebt.

Aus dem ersten Theil des Ausdrucks entsteht auf diese Weise von den nach Art.\ $2$ Vorschrift I gebildeten $\varepsilon$
\[
\begin{aligned}
&\Sigma\varepsilon\text{ wo }X\text{ ganz, }[Y]\text{ gerade}\\
&-\Sigma\varepsilon\text{ wo ausserdem }y\text{ ganz, }[x]\text{ gerade}
\end{aligned}
\]
Für die folgenden Theile kann
\[
\left.\begin{aligned}
&-(B)\mathrm{Int.}(p-B\mu\omega i)\text{ wenn }\xi=0\\
&+(B)\mathrm{Int.}(p+B\mu\omega i)\text{ wenn }\xi=\tfrac12
\end{aligned}\right\}\ X\text{ ganz, }[Y]\text{ gerade, }0<\eta<\tfrac12
\]
\[
\left.\begin{aligned}
&-(A)\mathrm{Int.}(p+A\mu\omega i)\text{ wenn }\eta=0\\
&+(A)\mathrm{Int.}(p-A\mu\omega i)\text{ wenn }\eta=\tfrac12
\end{aligned}\right\}\ X\text{ ganz, }[Y]\text{ gerade, }0<\xi<\tfrac12
\]
\[
\begin{gathered}
-\tfrac12[(A')+(B')]\mathrm{Int.}(p+A'\mu\omega i)\\
\text{wenn }X\text{ ganz, }[Y+A'\omega]\text{ gerade, }\xi\text{ und }\eta=0\text{ oder }\tfrac12
\end{gathered}
\]
gesetzt werden, worin $A'=+A$ oder $-A$ ist, wenn $\eta=0$ oder $\tfrac12$, $B'=+B$ oder $-B$, wenn $\xi=0$ oder $\tfrac12$, und worin z.\ B.\ $(A):+1$ oder $-1$ bezeichnet, jenachdem $A$ positiv oder negativ ist.

Multiplicirt man mit $(-1)^X$, führt die Summation über $X$ aus, lässt dabei in diesen Ausdrücken u.\ zwar im
\[
\begin{aligned}
\text{ersten}\quad&p-B\mu\omega i,\ P-BD\omega i,\ \pi-AB\omega i-BB\omega,\ X,\ [Y]\\
&\text{bez. in }ip,\ iP,\ i\pi,\ -Y,\ [X]\\[3pt]
\text{zweiten}\quad&p+B\mu\omega i,\ P+BD\omega i,\ \pi+AB\omega i+BB\omega,\ X,\ [Y]\\
&\dotstwo\quad p,\ P,\ \pi\quad X,\ [Y]\\[3pt]
\text{dritten}\quad&p+A\mu\omega i,\ P+AD\omega i,\ \pi+AA\omega i+AB\omega,\ X,\ [Y]\\
&\dotstwo\quad p,\ P,\ \pi\quad X,\ [Y]\\[3pt]
\text{vierten}\quad&p-A\mu\omega i,\ P-AD\omega i,\ \pi-AA\omega i-AB\omega,\ X,\ [Y]\\
&\dotstwo\text{ in }im-ip,\ iM-iP,\ i-i\pi,\ Y-B,\ A-1-[X]
\end{aligned}
\]

\SourceLeaf{389}%
übergehen und bezeichnet das aus dem fünften Ausdruck sich ergebende Resultat mit $Q_1$, so entsteht
\[
\begin{aligned}
&-\Sigma(-1)^Y(B)\mathrm{Int.}\,ip,\\
&\qquad\text{wo }Y\text{ ganz, }[X]\text{ gerade, }\eta=\omega,\ 0<\xi<\tfrac12\\
&+\Sigma(-1)^X(B)\mathrm{Int.}\,p,\\
&\qquad\text{wo }X\text{ ganz, }[Y]\text{ gerade, }\xi=\tfrac12+\omega,\ 0<\eta<\tfrac12\\
&-\Sigma(-1)^X(A)\mathrm{Int.}\,p,\\
&\qquad\text{wo }X\text{ ganz, }[Y]\text{ gerade, }\eta=\omega,\ 0<\xi<\tfrac12\\
&+\Sigma(-1)^Y(A)\mathrm{Int.}(im-ip),\\
&\qquad\text{wo }Y\text{ ganz, }[X]\text{ gerade, }\xi=\tfrac12+\omega,\ 0<\eta<\tfrac12\\
&+Q_1
\end{aligned}
\]

Die Untersuchung der einzelnen Fälle lässt erkennen, dass unter der Voraussetzung $M\equiv1\ (\text{mod. }2+2i)$
\[
\begin{aligned}
Q_1={}&-\mathrm{Int.}\,\mu\omega i\text{ ist, wenn }M\text{ im }1.\text{ Quadranten liegt}\\
&-\mathrm{Int.}(\tfrac12mi+\mu\omega i),\quad\text{wenn }M\text{ im }2.\text{ Quadr. und }\frac{M-1}{2}\text{ gerade}\\
&+\mathrm{Int.}(\tfrac12mi-\mu\omega i),\quad\text{wenn }M\text{ im }4.\text{ Quadr. und }\frac{M-1}{2}\text{ gerade}
\end{aligned}
\]
indem man eine complexe Zahl gerade oder ungerade nennt, je nachdem sie durch $2$ theilbar ist oder nicht.

Hiernach wird also bei Anwendung der in den Vorschriften II und IV bestimmten $\varepsilon$
\[
\begin{aligned}
\mathrm{Dec.}\frac mM={}&\mathrm{I},& &\Sigma\varepsilon,\ \text{wo }X\text{ ganz, }[Y]\text{ gerade}\\
&&&-4\Sigma\varepsilon,\ \text{wo noch }y\text{ ganz, }[x]\text{ gerade}\\
&\mathrm{II},&&-\Sigma\varepsilon\mathrm{Int.}\,ip,\ \text{wo }Y\text{ ganz, }[X]\text{ gerade}\\
&\mathrm{IV},&&+\Sigma\varepsilon\mathrm{Int.}\,p,\ \text{wo }X\text{ ganz, }[Y]\text{ gerade}\\
&\mathrm{II},&&+\Sigma\varepsilon\mathrm{Int.}\,p,\ \text{wo }X\text{ ganz, }[Y]\text{ gerade}\\
&\mathrm{IV},&&+\Sigma\varepsilon\mathrm{Int.}(im-ip),\ \text{wo }Y\text{ ganz, }[X]\text{ gerade}\\
&&&+Q_1
\end{aligned}
\]

In einer andern Form erhält man den Ausdruck für den Decidenten, wenn man zuerst nach $X$ summirt und dabei die Anzahl der zwischen $X'$ und $X''$ liegenden ungeraden Zahlen durch $[\tfrac12X''+\tfrac12]-[\tfrac12X'+\tfrac12]$ darstellt, nemlich
\[
\begin{aligned}
\mathrm{Dec.}\frac mM={}&\mathrm{I},& &\Sigma\varepsilon,\ \text{wo }Y\text{ ganz, }[X]\text{ ungerade}\\
&&&-4\Sigma\varepsilon,\ \text{wo noch }y\text{ ganz, }[x]\text{ gerade}\\
&\mathrm{II},&&-\Sigma\varepsilon\mathrm{Int.}\,ip,\ \text{wo }X\text{ ganz, }[Y]\text{ gerade}\\
&\mathrm{IV},&&+\Sigma\varepsilon\mathrm{Int.}\,p,\ \text{wo }Y\text{ ganz, }[X]\text{ ungerade}\\
&\mathrm{II},&&+\Sigma\varepsilon\mathrm{Int.}\,p,\ \text{wo }Y\text{ ganz, }[X]\text{ ungerade}\\
&\mathrm{IV},&&+\Sigma\varepsilon\mathrm{Int.}(im-ip),\ \text{wo }X\text{ ganz, }[Y]\text{ ungerade}\\
&&&+Q_2
\end{aligned}
\]
\[
\begin{aligned}
Q_2={}&-\mathrm{Int.}(-\mu\omega),\quad\text{wenn }M\text{ im }2.\text{ Quadr.}\\
&+\mathrm{Int.}(\tfrac12m-\mu\omega),\quad\text{wenn }M\text{ im }1.\text{ Quadr. und }\frac{M-1}{2}\text{ ungerade}\\
&-\mathrm{Int.}(\tfrac12m+\mu\omega),\quad\text{wenn }M\text{ im }3.\text{ Quadr. und }\frac{M-1}{2}\text{ ungerade}
\end{aligned}
\]

\SourceLeaf{390}%
Führt man die Summation nach $Y$ zuerst aus, wählt aber die zweite so eben angewandte Art der Bestimmung der Anzahl der zwischen zwei Werthen liegenden ungeraden Zahlen, so wird
\[
\begin{aligned}
\mathrm{Dec.}\frac mM={}&\mathrm{I},& &\Sigma\varepsilon,\ \text{wo }X\text{ ganz, }[Y]\text{ ungerade}\\
&&&-4\Sigma\varepsilon,\ \text{wo noch }y\text{ ganz, }[x]\text{ gerade}\\
&\mathrm{II},&&-\Sigma\varepsilon\mathrm{Int.}\,ip,\ \text{wo }Y\text{ ganz, }[X]\text{ ungerade}\\
&\mathrm{IV},&&+\Sigma\varepsilon\mathrm{Int.}\,p,\ \text{wo }X\text{ ganz, }[Y]\text{ ungerade}\\
&\mathrm{II},&&+\Sigma\varepsilon\mathrm{Int.}\,p,\ \text{wo }X\text{ ganz, }[Y]\text{ ungerade}\\
&\mathrm{IV},&&+\Sigma\varepsilon\mathrm{Int.}(im-ip),\ \text{wo }Y\text{ ganz, }[X]\text{ ungerade}\\
&&&+Q_3
\end{aligned}
\]
\[
\begin{aligned}
Q_3={}&-\mathrm{Int.}(-\mu\omega i),\quad\text{wenn }M\text{ im }3.\text{ Quadr.}\\
&-\mathrm{Int.}(\tfrac12mi+\mu\omega i),\quad\text{wenn }M\text{ im }2.\text{ Quadr. und }\frac{M-1}{2}\text{ ungerade}\\
&+\mathrm{Int.}(\tfrac12mi-\mu\omega i),\quad\text{wenn }M\text{ im }4.\text{ Quadr. und }\frac{M-1}{2}\text{ ungerade}
\end{aligned}
\]

Summirt man zuerst noch $X$ und gebraucht dabei die erste Art der Darstellung der Anzahl der zwischen zwei Werthen liegenden ungeraden Zahlen, so erhält man die in [V.] Art.\ $15$ angegebene Form für den Decidenten, wo die Grösse $q$ auch durch folgende Gleichung definirt werden kann
\[
\begin{aligned}
q={}&-\mathrm{Int.}\,\mu\omega,\quad\text{wenn }M\text{ im }4.\text{ Quadr.}\\
&+\mathrm{Int.}(\tfrac12m-\mu\omega),\quad\text{wenn }M\text{ im }1.\text{ Quadr. und }\frac{M-1}{2}\text{ gerade}\\
&-\mathrm{Int.}(\tfrac12m+\mu\omega),\quad\text{wenn }M\text{ im }3.\text{ Quadr. und }\frac{M-1}{2}\text{ gerade}
\end{aligned}
\]

Die Vereinigung dieser vier Ausdrücke für den Decidenten bildet das in [V.] Art.\ $7$ aufgestellte Resultat, weil $\mathrm{Int.}\,ip-\mathrm{Int.}\,p$ gleich $3$ wird für $[x]$ gerade $[y]$ ungerade, sonst aber gleich $1$, ferner $\mathrm{Int.}(im-ip)+\mathrm{Int.}\,p$ gleich $0$ für $[x]$ gerade $[y]$ gerade, in den übrigen Fällen aber gleich $4$.

[V.] Art.\ $[7.]$ Die erste Tafel für das Beispiel giebt in der ersten Spalte die zu jedem ganzzahligen $P$ zugehörigen Werthe von $\frac{37\pdot P}{M}$ oder $37(\xi+\eta i)$, wenn $0<\xi<\tfrac12$, $0<\eta<\tfrac12$ ist, in der zweiten $\frac{37\pdot Pm}{M}$ oder $37\pdot p$, in der dritten die in $p$ enthaltene grösste ganze Zahl, in der vierten $\pm\mathrm{Int.}\,p$, wo das obere Zeichen gilt, wenn $P$ durch $1+i$ theilbar, das untere, wenn $P$ nicht durch $1+i$ theilbar ist.

[V.] Art.\ $[9.]$. $[12.]$ Die verbesserte Bezeichnungsweise der $\theta$ ist nur bei der zweiten und dritten Classe Artt.\ $9$. $10$ angedeutet, aber auch auf die erste und vierte Artt.\ $8$. $11$ auszudehnen. Hiernach wird ein $\theta^\lambda=T+Ui$ denjenigen Index $\lambda,=0,1,2$ oder $3$ haben, für welchen die durch die Gleichungen

\SourceLeaf{391}%
\[
\begin{aligned}
i^\lambda M&=\mathfrak A+\mathfrak B i,\quad i^{-\lambda}\mu=\rho+\sigma i\\
\sigma\varphi^0&=-\text{Coëff. Img }\theta^\lambda(a-bi)=+bT-aU\\
\sigma\Phi^0&=+\text{Coëff. Img }\theta^\lambda(\mathfrak A-\mathfrak B i)=-\mathfrak B T+\mathfrak A U
\end{aligned}
\]
bestimmten Grössen $\varphi^0$ und $\Phi^0$ zwischen $0$ und $\tfrac12$ liegen.

Um nach den Andeutungen in Art.\ $9$ $(3)$ zu beweisen, dass, wenn $T$, $U$ zwei ganze reelle Zahlen sind, welche die so eben aufgestellten Bedingungen erfüllen, $T+Ui$ sich auch in dem bei einer der vier Combinationen Artt.\ $8$. $11$. bestimmten Complexus $\theta^\lambda$ befindet, bezeichne man mit $\varphi'$, $\Phi'$ diejenigen ganzen complexen Zahlen, für welche die Gleichung
\[
T+Ui=\varphi'i^\lambda M+\Phi'm
\]
Statt hat und für welche eine der vier Grössen $\pm\frac{\varphi'-\varphi^0}{m}$, $\pm i\frac{\varphi'-\varphi^0}{m}$ so beschaffen, dass der reelle Theil und der Coëfficient des imaginären Theils zwischen $0$ und $\tfrac12$ liegen (\textit{Theoria residuorum biquadr. artt.}\ $45,46$). Die betreffende Grösse ist dann, wie man aus der Untersuchung der in den vier Combinationen enthaltenen sechzehn einzelnen Fälle leicht ersieht, $\frac pm$ und die ihr entsprechende Grösse unter $\pm\frac{\Phi'-\Phi^0}{m}$, $\pm i\frac{\Phi'-\Phi^0}{M}$ ist $\frac PM$, weil $\frac{\Phi'-\Phi^0}{M}=-\frac{\varphi'-\varphi^0}{m}i^\lambda$ wird.

Aus dieser Art der Darstellung der Grössen $\frac pm$ oder $\frac PM$ folgt auch, dass $\mathrm I,\ \Sigma\varepsilon$ von allen aus $\Theta+\Theta'+\Theta''+\Theta'''$ besteht, worin $\Theta^\lambda$ die Summe derjenigen $\varepsilon$ bedeutet, die für jeden Ganzepunkt $\theta$ innerhalb des Parallelogramms $0,\ \tfrac12m,\ \tfrac12m+\tfrac12i^\lambda M,\ \tfrac12i^\lambda M,$
\begin{itemize}
\item[] $=+1$ zu setzen sind, wenn $\theta$ durch $1+i$ theilbar und $\text{Coëff. Imag. }\mu i^{-\lambda}$ positiv oder wenn keine Bedingung gilt, dagegen
\item[] $=-1$ wenn nur eine gilt.
\end{itemize}

[V.] Art.\ $13$. Die Bestimmung von $\Theta^\lambda$ kann entweder durch die oben für $\mathrm{Dec.}\frac mM$ angewandten vier verschiedenen Summationsarten oder, was im Wesentlichen dasselbe ist, nach den in [II.] Art.\ $11$ angedeuteten Methoden ausgeführt werden, bei welchen dann die vier Constructionen zu Grunde zu legen sind, die durch Verbindung der Punkte, deren $\theta$ ein Vielfaches von $1+i$ ist, resp.\ mit den Punkten $\theta+1$, $\theta+i$, $\theta-1$ und $\theta-i$ entstehen.

Lässt man in der Begrenzung des zuvor erwähnten Parallelogramms allen den Punkten ein $\theta$ entsprechen, für welche der reelle oder imaginäre Theil von $\theta$ eine ganze Zahl wird, bezeichnet mit $\theta^0$ die nächste durch $1+i$ theilbare Ganze bei $\theta$, mit $l$ die Ortsverschiebung von einem Punkte des geraden Begrenzungsstückes, das den Punkt $\theta$ enthält, bis zu irgend einem nachfolgenden Punkte derselben Geraden, also z.\ B.\ bei jenem Parallelogramm der Reihe nach die Grössen $m,\ Mi^\lambda,\ -m,\ -Mi^\lambda$, und setzt
\[
\varepsilon=\pm1\text{ mit dem Zeichen des imaginären Theils von }\frac l{\theta-\theta^0}
\]
so ergiebt die Vereinigung der auf die eine oder andere Weise erhaltenen vier Resultate $4\Theta^\lambda=-\Sigma\varepsilon$.

Die gesonderte Bestimmung der den Eckpunkten entsprechenden $\theta$ und $\varepsilon$ wird umgangen, wenn man dies Parallelogramm durch ein anderes ersetzt, dessen Begrenzungen den Begrenzungen des erstern
\SourceLeaf{392}%
unendlich nahe sind, und welches die beiden Punkte $0$ und $\tfrac12m+\tfrac12i^\lambda M$ nicht einschliesst. Die Begrenzung eines solchen Parallelogramms erhält man, wenn man sie an die positiven Seiten der Linien
\[
\begin{gathered}
0\dotsfour\tfrac12m,\quad\tfrac12m+\tfrac12i^\lambda M\dotsfour\tfrac12m,\\
\tfrac12m+\tfrac12i^\lambda M\dotsfour\tfrac12i^\lambda M,\quad 0\dotsfour\tfrac12i^\lambda M
\end{gathered}
\]
legt. Lässt man den vier so entstandenen Geraden der Reihe nach die unendlich kleinen positiven Grössen $\omega_1,\omega_2,\omega_3,\omega_4$ entsprechen, so kann man für die auf ihnen liegenden Punkte $\theta$
\[
\begin{aligned}
\theta&=p=m(\xi+\omega_1i),\\
-\theta i^\lambda+\tfrac12mi^\lambda+\tfrac12M&=P=M(\xi+\omega_2i),\\
-\theta+\tfrac12m+\tfrac12Mi^\lambda&=p=m(\xi+\omega_3i),\\
\theta i^{-\lambda}&=P=M(\xi+\omega_4i)\quad\text{wenn }\lambda\text{ gerade}
\end{aligned}
\]
\[
\begin{aligned}
\theta&=p=m(\xi+\omega_1i),\\
\theta i^{1-\lambda}-\tfrac12mi^{1-\lambda}+\tfrac12M&=P=M(\tfrac12+\omega_2+\eta i),\\
\theta i+\tfrac12m-\tfrac12Mi^{1+\lambda}&=p=m(\tfrac12+\omega_3+\eta i),\\
\theta i^{-\lambda}&=P=M(\xi+\omega_4i)\quad\text{wenn }\lambda\text{ ungerade}
\end{aligned}
\]
setzen, worin $\xi$ und $\eta$ auch theilweise zur Schliessung der Figur das Gebiet der reellen Werthe von $0$ bis $\tfrac12$ um unendlich kleine Grössen überschreiten.

Bezeichnen $G$, $g$, $H$, $h$ die Summen der resp.\ nach den Vorschriften II, III, IV, V (in Artt.\ $3$ bis $6$) gebildeten $\varepsilon$, und umfassen $G'$ oder $G_{\prime}$ und $g'$ oder $g_{\prime}$ diejenigen $\varepsilon$, welche für die beim zweiten Parallelogramm etwa auftretenden unendlich kleinen Werthe von $\xi$ Statt haben, im Uebrigen aber resp.\ nach den Vorschriften II und III gebildet sind, beziehen sich ferner $G''$ oder $G_{\prime\prime}$ und $g''$ oder $g_{\prime\prime}$ ebenso auf dieselben Vorschriften aber auf die unendlich kleinen Werthe von $\tfrac12-\xi$, und endlich $H',h_{\prime},H'',h_{\prime\prime}$ resp.\ auf die Vorschriften IV, V, IV, V und die unendlich kleinen Werthe resp.\ von $\tfrac12-\eta,\tfrac12-\eta,\eta,\eta$, so wird
\[
\begin{aligned}
4\Theta^\lambda={}&-(g+g'+g'')-i^\lambda(G+G'+G'')\\
&+i^\lambda(g+g_{\prime}+g_{\prime\prime})+(G+G_{\prime}+G_{\prime\prime})\quad\text{wenn }\lambda\text{ gerade}\\
4\Theta^\lambda={}&-(g+g'+g'')-i^{1-\lambda}(H+H'+H'')\\
&-i^{\lambda-1}(h+h_{\prime}+h_{\prime\prime})+(G+G_{\prime}+G_{\prime\prime})\quad\text{wenn }\lambda\text{ ungerade}
\end{aligned}
\]

Für denjenigen Eckpunkt $\theta$ des Parallelogramms, welcher dem Punkte $0$ zunächst liegt, bezeichne $\xi_1$ den zugehörigen Werth von dem $\xi$ der ersten Seite, $\xi_4$ den zugehörigen Werth von dem $\xi$ der vierten Seite, so dass
\[
\theta=m(\xi_1+\omega_1i)=i^\lambda M(\xi_4+\omega_4i)
\]
wird, dann ergiebt sich dasjenige $\xi$, welchem auf der ersten Seite oder deren Verlängerung ein Punkt $p$ mit dem reellen Theile gleich $0$ entspricht, aus der Gleichung
\[
(\mathrm{Real.}\,p=0),\quad\xi-\xi_1=\sigma a\mathfrak A\omega_1-\sigma\omega_4
\]
worin die positiven Factoren der unendlich kleinen positiven Grössen durch die Einheit ersetzt sind und $\sigma$, $\mathfrak A$ die durch
\[
\rho+\sigma i=i^{-\lambda}(\alpha+\beta i),\quad\mathfrak A+\mathfrak B i=i^\lambda(A+Bi)
\]
bestimmten reellen Grössen bedeuten. Dieser Punkt $p$ liegt auf der ersten Seite selbst, wenn $\xi-\xi_1$ positiv, also, indem man $\omega_1$ unendlich klein gegen $\omega_4$ annimmt, wenn $\sigma$ negativ ist. Der dem Punkte $p$ zunächst liegende Punkt $p^0$, dessen darstellende Zahl durch $1+i$ getheilt wird, ist der Punkt $0$, also hat
\SourceLeaf{393}%
$\mathrm{Imag.}\frac m{p-p^0}$ oder $\mathrm{Imag.}\frac1{\xi+\omega_1i}$ das Minuszeichen. Man erhält daher für $\mathrm{Real.}\,p=0$:
\[
\varepsilon=-1\text{ wenn }(\sigma)=-1,\quad\varepsilon=0\text{ wenn }(\sigma)=+1,\quad\text{d. i. }\varepsilon=-\tfrac12+\tfrac12(\sigma)
\]
und auf dieselbe Weise für $\mathrm{Imag.}\,p=0$
\[
\begin{gathered}
\xi-\xi_1=\sigma b\mathfrak B\omega_1-\sigma\omega_4,\quad\mathrm{Imag.}\frac m{p-p^0}=\mathrm{Imag.}\frac1{\xi+\omega_1i},\\
\varepsilon=-\tfrac12+\tfrac12(\sigma)\quad\text{also }g'=-1+(\sigma)
\end{gathered}
\]

In Bezug auf die vierte Seite wird $P^0=0$, $\mathrm{Imag.}\frac M{P-P^0}=\mathrm{Imag.}\frac1{\xi+\omega_4i}$
\[
\begin{aligned}
&\text{also für }\mathrm{Real.}(i^\lambda P)=0;&&\xi-\xi_4=-\sigma a\mathfrak A\omega_4+\sigma\omega_1,\\
&&&\varepsilon=-\tfrac12+\tfrac12(\sigma a\mathfrak A)\\
&\text{und für }\mathrm{Imag.}(i^\lambda P)=0;&&\xi-\xi_4=-\sigma b\mathfrak B\omega_4+\sigma\omega_1,\\
&&&\varepsilon=-\tfrac12+\tfrac12(\sigma b\mathfrak B)
\end{aligned}
\]
demnach $G_{\prime}=-1+\tfrac12(\sigma a\mathfrak A)+\tfrac12(\sigma b\mathfrak B)$ oder, weil $\rho=a\mathfrak A+b\mathfrak B$ ist, $G_{\prime}=-1+\tfrac12(\rho\sigma)+\tfrac12(\rho\sigma ab\mathfrak A\mathfrak B)$

Der Theil $R_1^\lambda$ von $4\Theta^\lambda$, der aus dem unendlich nahe bei dem Punkte $0$ liegenden Stücke der Begrenzung entsteht, ist also
\[
R_1^\lambda=+G_{\prime}-g'=-(\sigma)+\tfrac12(\rho\sigma)+\tfrac12(\rho\sigma ab\mathfrak A\mathfrak B)
\]

Durch ähnliche Betrachtungen findet man für die Theile $R_2^\lambda$, $R_3^\lambda$, $R_4^\lambda$, welche ebensolche Beziehungen resp.\ zu den Punkten $\tfrac12m,\tfrac12m+\tfrac12Mi^\lambda,\tfrac12Mi^\lambda$ haben, wie $R_1^\lambda$ zum Punkte $0$, bei geradem $\lambda$
\[
\begin{aligned}
R_2^\lambda&=-g''-i^\lambda G''=-\tfrac12(b)-\tfrac12(\mathfrak B),\\
R_3^\lambda&=-i^\lambda G'+i^\lambda g_{\prime}=i^\lambda(\sigma)-\tfrac12i^\lambda(\rho\sigma)-\tfrac12i^\lambda(\rho\sigma ab\mathfrak A\mathfrak B)\\
R_4^\lambda&=+i^\lambda g_{\prime\prime}+G_{\prime\prime}=\tfrac12i^\lambda(b)+\tfrac12i^\lambda(\mathfrak B)
\end{aligned}
\]
bei ungeradem $\lambda$
\[
\begin{aligned}
R_2^\lambda&=-g''-i^{1-\lambda}H''=-\tfrac12(b)-\tfrac12(\mathfrak B),\\
R_3^\lambda&=-i^{1-\lambda}H'-i^{\lambda-1}h_{\prime}=0,\\
R_4^\lambda&=-i^{\lambda-1}h_{\prime\prime}+G_{\prime\prime}=\tfrac12i^{\lambda+1}(a)+\tfrac12i^{\lambda+1}(\mathfrak A)
\end{aligned}
\]

[V.] Art.\ $[14.]$ Die Auswerthung der Summen von den nach Vorschrift III gebildeten $\varepsilon$ ergiebt sich aus der durch die Definition der $\varepsilon$ leicht zu verificirenden Gleichung
\[
\begin{aligned}
&\mathrm{III},\quad\Sigma\varepsilon\text{ von allen }-4\Sigma\varepsilon\text{ von denen, wo }y\text{ ganz, }[x]\text{ gerade}\\
&=\mathrm{III},\quad\Sigma[-\mathrm{Int.}(p-m\omega)+\mathrm{Int.}(p+m\omega)]
\end{aligned}
\]
worin $p$ alle Werthe annimmt, die den unter Vorschrift III angegebenen Bedingungen genügen. Diese Intensoren lassen sich nemlich mit Ausnahme der beiden dem kleinsten $(\xi^{*})$ und dem grössten zulässigen Werthe $(\xi^{**})$ von $\xi$ entsprechenden Intensoren, welche resp.\ gleich
\[
\begin{gathered}
-\mathrm{Int.}(p^{*}-m\omega)\text{ und }+\mathrm{Int.}(p^{**}+m\omega)\text{ oder}\\
-\mathrm{Int.}(\omega m)\text{ und }+\mathrm{Int.}(\tfrac12-\omega)m
\end{gathered}
\]
sind, immer zu je zweien $+\mathrm{Int.}(p'+m\omega)$ und $-\mathrm{Int.}(p''-m\omega)$ so zusammen ordnen, dass zwischen $\xi'$ und $\xi''$, welche den Grössen $p'$ und $p''$ entsprechen, kein Werth von $\xi$ liegt, der den reellen oder imaginären Theil von $p$ zu einer ganzen Zahl macht, so dass also die zwei Intensoren sich stets gegenseitig annulliren.

\SourceLeaf{394}%
[V.] Art.\ $[15.]$ Es ist
\[
\begin{aligned}
\mathrm{II},\quad&+\Sigma\varepsilon\mathrm{Int.}\,ip\text{ wo }Y\text{ ganz }[X]\text{ gerade},\\
&-\Sigma\varepsilon\mathrm{Int.}\,ip\text{ wo }X\text{ ganz }[Y]\text{ ungerade}\\
={}&\Sigma[-\mathrm{Int.}\,i(p-m\omega)+\mathrm{Int.}\,i(p+m\omega)]\\
&\text{für diejenigen }p,\text{ für welche }x\text{ oder }y\text{ ganz, }[X]\text{ gerade}\\
&Y\text{ ungerade, }0<\xi<\tfrac12,\quad\eta=\omega\\
&+\mathrm{Int.}\,i(p^{*}-m\omega)\text{ wenn }[X^{*}]\text{ gerade }[Y^{*}]\text{ ungerade}\\
&-\mathrm{Int.}\,i(p^{**}+m\omega)\text{ wenn }[X^{**}]\text{ gerade }[Y^{**}]\text{ ungerade}
\end{aligned}
\]
wie man sich leicht überzeugt, wenn man auf der zweiten Seite der Gleichung die Summation nach dem in der vorhergehenden Note angewandten Verfahren über jedes so kleine Intervall von $\xi_{\prime}$ bis $\xi_{\prime\prime}$ ausführt, dass es zwischen $\xi_{\prime}$ und $\xi_{\prime\prime}$ kein $\xi$ gibt, welches in dem zugehörigen $P$ den reellen oder imaginären Theil zu einer ganzen Zahl macht. Die Anwendung der nach Vorschrift III gebildeten $\varepsilon$ lässt die zweite Seite dieser Gleichung die in Art.\ $15$. aufgestellte Form annehmen.

[V.] Art.\ $[15.]$ Die Verwandlung der Summen von den nach Vorschrift IV gebildeten $\lambda'\varepsilon$ in die Summen der $\varepsilon$ aus V ergiebt sich durch eben solche Betrachtungen wie die in der letzten Note angewandten, wenn noch die Gleichung
\[
\begin{aligned}
&\mathrm V,\quad\Sigma\varepsilon\text{ von allen, }-4\Sigma\varepsilon\text{ wo }x\text{ ganz }[y]\text{ gerade},\\
&=+\mathrm{Int.}(\tfrac12im+m\omega)-\mathrm{Int.}(\tfrac12im+\tfrac12-\omega m)
\end{aligned}
\]
zu Hülfe gezogen wird, die der zuvor ermittelten Auswerthung der Summe von den $\varepsilon$ in Vorschrift III entspricht.

[V.] Art.\ $[17.]$ Bestimmt man die Hülfsgrössen $U$, $T$, $L$, $V$ durch die Gleichungen
\[
\begin{aligned}
U={}&1-2(B)-(AB)+(\alpha a)+(\beta b)+(\alpha b)-(\beta a)\\
&+(\alpha\beta ab)-(\alpha\beta abAB)\quad\text{oder}\\
U={}&(1+(A))(1-(B))(1-(a\beta)+(b\beta)+(\alpha\beta))
\end{aligned}
\]
\[
\begin{gathered}
\text{weil }(\alpha a)+(\beta b)=(A)+(A\alpha\beta ab),\\
(\alpha b)-(\beta a)=(B)-(B\alpha\beta ab)\text{ ist},
\end{gathered}
\]
\[
\begin{aligned}
T&=-2+(a)+(b)-(\beta)-2(B)-(\alpha ab)\quad\text{oder}\\
T&=-2+(a)+(b)-(\beta)-2(B)-(aA)+(bB)-(\beta AB)\\
L&=(1+(A))(1+(B))(1+(\beta))-(1-(A))(1-(B))(1-(\beta))\\
V&=(1+(A))(1-(B))(-2(a)+2(\beta)+(ab)+(\alpha\beta))\quad\text{oder}\\
V&=(1+(A))(1-(B))(-(a)+(b)+(\beta)+(ab)-(\beta ab)+(\alpha\beta))
\end{aligned}
\]
\[
\text{weil }(a)+(b)=(\beta)+(\beta ab)\text{ wenn }A\text{ positiv }B\text{ negativ}
\]
und bezeichnet mit $W'$, $S'$, $Q'$ die Grössen, in welche die $W$, $S$, $Q$ des Ausdrucks für den $\mathrm{Dec.}\frac mM$ in Art.\ $14$. übergehen, wenn man darin $m$ mit $M$ also $\alpha+\beta i$ mit $\alpha-\beta i$ vertauscht, so wird

\SourceLeaf{395}%
\[
\begin{aligned}
2W'&=-5(B)-(AB)\quad\text{wenn }\frac{M-1}{2}\text{ gerade}\\
2W'&=-(B)-(AB)\quad\text{wenn }\frac{M-1}{2}\text{ ungerade}\\
2S'&=-(\alpha\beta abAB)-(\beta)-(B)+(b)\\
2Q'+2S'+2W'&=2T+U\quad\text{wenn }\frac{M-1}{2}\text{ gerade}\\
2Q'+2S'+2W'&=-4+4(a)+U\quad\text{wenn }\frac{M-1}{2}\text{ ungerade}\\
8\psi&=8(q+r+s+w)-(2Q'+2S'+2W')
\end{aligned}
\]

Ersetzt man hier $8(q+r+s+w)$ durch dessen in Art.\ $15$ aufgestellten Werth, bringt ihn aber unter die Form
\[
\begin{array}{lll}
2T+4L+V&\text{wenn }\dfrac{M-1}{2}\text{ gerade}&\dfrac{m-1}{2}\text{ gerade}\\[4pt]
2T-16+16(A)+V&\text{wenn }\dfrac{M-1}{2}\text{ gerade}&\dfrac{m-1}{2}\text{ ungerade}\\[4pt]
-4+4(a)+V&\text{wenn }\dfrac{M-1}{2}\text{ ungerade}&\dfrac{m-1}{2}\text{ gerade}\\[4pt]
-20+4(a)+V&\text{wenn }\dfrac{M-1}{2}\text{ ungerade}&\dfrac{m-1}{2}\text{ ungerade}
\end{array}
\]
und beachtet, dass
\[
V-U=-\tfrac12(1+(a))(1+(A))(1+(\alpha))(1-(b))(1-(B))(1-(\beta))
\]
ist, so erhält man für $\psi$ die in Art.\ $17$ angegebene Bestimmungsart.

[VI.] Art.\ $3$. Das unvollständige Citat kann auf Art.\ $4$ des Bruchstücks III bezogen werden.
\begin{flushright}\textsc{Schering.}\end{flushright}
\BookRule{.3\textwidth}

